\subsection{章末確認問題}

\begin{enumerate}
  \item ソフトウェアエンジニアリング運用は、単なる本番配備とどう違うか。
  \item 運用計画、運用提供、運用制御の違いを説明せよ。
  \item 運用担当技術者とソフトウェア技術者の役割の違いを説明せよ。
  \item IaCが環境差分を減らす理由を説明せよ。
  \item 配備とリリースの違いを、具体例を使って説明せよ。
  \item カナリアリリースがリスクを減らす理由を説明せよ。
  \item ロールバック計画にデータ移行が含まれるべき理由を説明せよ。
  \item インシデント管理、問題解決、変更管理の関係を説明せよ。
  \item 監視とテレメトリが、インシデント予防に役立つ理由を説明せよ。
  \item 小規模組織では、どの運用活動から始めるのが現実的かを述べよ。
\end{enumerate}

\begin{pointbox}{解答の方向性}
1は、配備、設定、監視、障害対応、変更管理、容量管理、バックアップ、復旧、セキュリティを含む点を述べる。2は、計画が準備、提供が実行、制御が監視と管理である点を述べる。5は、配備は環境へ置くこと、リリースは利用者に使わせることと説明する。8は、インシデントを記録し、問題として原因を分析し、変更管理で修正を安全に入れる流れとして説明できる。
\end{pointbox}
